%% kapa: exam=1 round=24 year=2013 period=1 form=A subject=economics subject_ko=경제원론 seq=27 number=67
%% source: Q-Net 감정평가사 기출문제 (cst00309) artlSeq=5073802 file=extracted/1008/5073802/2013년 제24회 감정평가사 1차시험 기출문제/1교시(A형).hwp
%% conversion: kordoc 4.24.51 -> tools/gen1.py (auto); review=proofread-sonnet; flags=none
\begin{question}{67}
유동성 함정에 관한 설명으로 옳지 않은 것은?
\begin{choices}
  \item 화폐수요의 이자율 탄력성이 무한대인 경우에 발생한다.
  \item 채권의 가격이 매우 높아서 추가적인 통화공급이 투기적 화폐수요로 모두 흡수된다.
  \item 이자율이 매우 낮아 향후 이자율이 상승할 것으로 예상될 경우 유동성 함정이 발생할 수 있다.
  \item 확장적 통화정책은 이자율을 하락시키지 못하여 총수요 확대효과가 없다.
  \item 확장적 재정정책은 이자율을 상승시켜 총수요 확대효과가 없다.
\end{choices}
\end{question}
\answer{5}
